\section*{第 \ref{ch:b1:potential-capacitors} 章——电势；导体与电容器}

\begin{solution}{exo:b1:potential-capacitors:1}
$V = 8.99 \times 10^9 \times 10^{-6}/1 = \qty{9.0}{kV}$; $W = qV = \qty{9.0}{mJ}$;
释放后，$\tfrac12mv^2 = \qty{9}{mJ}$：$v = \sqrt{18} = \qty{4.2}{m/s}$。
\end{solution}

\begin{solution}{exo:b1:potential-capacitors:2}
$E = 100/10^{-3} = \qty{1e5}{V/m}$; $E_k = eU = \qty{100}{eV} = \qty{1.6e-17}{J}$,
故 $v = \sqrt{2E_k/m_e} = \qty{5.9e6}{m/s}$；做匀加速运动，
$t = 2d/v = \qty{3.4e-10}{s}$.
\end{solution}

\begin{solution}{exo:b1:potential-capacitors:3}
$E_A = 10/0.002 = \qty{5}{kV/m}$，$E_B = 10/0.008 = \qty{1.25}{kV/m}$。电场
指向电势降低的方向；电子（$q < 0$）
向相反方向加速，朝 $V$ 增大的方向。
\end{solution}

\begin{solution}{exo:b1:potential-capacitors:4}
$C = 4\pi\varepsilon_0R = 1.11 \times 10^{-10} \times 0.15 = \qty{17}{pF}$; $Q = CV =
\qty{3.3}{\micro C}$；$E = V/R = 2 \times 10^5/0.15 = \qty{1.3}{MV/m}$。击穿
发生于 $V = E_{\max}R = 3 \times 10^6 \times 0.15 = \qty{450}{kV}$ 时——大型装置
采用大球。
\end{solution}

\begin{solution}{exo:b1:potential-capacitors:5}
圆环上每一元素到该点的距离均为 $\sqrt{R^2 + z^2}$：$V =
Q/4\pi\varepsilon_0\sqrt{R^2 + z^2}$; $E_z = -\dd V/\dd z = Qz/4\pi\varepsilon_0(R^2 +
z^2)^{3/2}$。圆盘：$\dd q = 2\pi\sigma a\,\dd a$ 的电荷位于距离 $\sqrt{a^2 + z^2}$ 处：$V =
\dfrac{\sigma}{2\varepsilon_0}\displaystyle\int_0^R\dfrac{a\,\dd a}{\sqrt{a^2 + z^2}} =
\dfrac{\sigma}{2\varepsilon_0}\bigl(\sqrt{z^2 + R^2} - z\bigr)$，且 $-\dd V/\dd z =
\dfrac{\sigma}{2\varepsilon_0}\Bigl(1 - \dfrac{z}{\sqrt{z^2 + R^2}}\Bigr)$.
\end{solution}

\begin{solution}{exo:b1:potential-capacitors:6}
$V(r) = V(R) + \displaystyle\int_r^R\dfrac{Qr'\,\dd r'}{4\pi\varepsilon_0R^3} =
\dfrac{Q}{4\pi\varepsilon_0R} + \dfrac{Q(R^2 - r^2)}{8\pi\varepsilon_0R^3} = \dfrac{Q}{4\pi\varepsilon_0}
\dfrac{3R^2 - r^2}{2R^3}$。金核：$V(0) = 1.5 \times 8.99 \times 10^9 \times 1.26 \times
10^{-17}/7 \times 10^{-15} = \qty{24}{MV}$：质子到达中心需 \qty{24}{MeV}，α 粒子
（$2e$）到达中心需 \qty{48}{MeV}，触及表面需 \qty{32}{MeV}。
卢瑟福的 \qty{5}{MeV} α 粒子在 $r$ 处折返，此处 $2eV =
\qty{5}{MeV}$：$V = \qty{2.5}{MV}$，$r = \qty{45}{fm}$——远在
原子核之外，这正是他的点电荷分析得以奏效的原因。
\end{solution}

\begin{solution}{exo:b1:potential-capacitors:7}
轴线上：$2p/4\pi\varepsilon_0r^3 = 2 \times 8.99 \times 10^9 \times 6.2 \times 10^{-30}/10^{-27} =
\qty{1.1e8}{V/m}$；垂直方向上减半，为 \qty{5.6e7}{V/m}。力矩 $pE =
\qty{6.2e-24}{N.m}$；能量差 $2pE = \qty{1.2e-23}{J}$ 与 $k_BT =
\qty{4.1e-21}{J}$：小 $330$ 倍——净取向仅 $0.1\%$。
位于 \qty{0.3}{nm} 处的钠离子：$E = e/4\pi\varepsilon_0r^2 = \qty{1.6e10}{V/m}$，$-pE
= \qty{-1.0e-19}{J} = \qty{-0.6}{eV}$，$24\,k_BT$：水分子
锁定到离子上——即水合作用。
\end{solution}

\begin{solution}{exo:b1:potential-capacitors:8}
$Q = 4\pi\varepsilon_0R^2E = 10^{-4} \times 3 \times 10^6/8.99 \times 10^9 = \qty{33}{nC}$,
$V = ER = \qty{30}{kV}$。两球连接后：$E_1 = V/R_1 = \qty{0.2}{MV/m}$，
$E_2 = V/R_2 = \qty{20}{MV/m} > \qty{3}{MV/m}$：小球使
空气电离并泄漏电荷——即\omterm{ex:b1:potential-capacitors:point}{尖端效应}。
\end{solution}

\begin{solution}{exo:b1:potential-capacitors:9}
$C = \varepsilon_r\varepsilon_0S/d = 2.2 \times 8.85 \times 10^{-12}/10^{-5} = \qty{2.0}{\micro F}$;
卷绕只改变形状，不改变正对面积：$C$ 不变（若加入第二层薄膜，
使每片箔两侧均与另一箔片正对，则\omterm{def:b1:potential-capacitors:capacitance}{电容}加倍）。
$W = \tfrac12 \times 2 \times 10^{-6} \times 400^2 = \qty{0.16}{J}$。除颤器：
$Q = 150 \times 10^{-6} \times 2000 = \qty{0.30}{C}$, $W = \tfrac12CU^2 = \qty{300}{J}$,
$P = 300/0.005 = \qty{60}{kW}$.
\end{solution}

\begin{solution}{exo:b1:potential-capacitors:10}
(a) $E = Q/4\pi\varepsilon_0r^2$，$U = \dfrac{Q}{4\pi\varepsilon_0}\Bigl(\dfrac1{R_1} -
\dfrac1{R_2}\Bigr)$，$C = 4\pi\varepsilon_0R_1R_2/(R_2 - R_1)$；当 $R_2 - R_1 = d \ll
R_1$，$R_1R_2 \approx R^2$，且 $C \approx 4\pi R^2\varepsilon_0/d = \varepsilon_0S/d$；当 $R_2 \to
\infty$，$C \to 4\pi\varepsilon_0R_1$。(b) $C/L = 2\pi\varepsilon_r\varepsilon_0/\ln(b/a) = 2\pi \times 2.3
\times 8.85 \times 10^{-12}/\ln5 = \qty{79}{pF/m}$，量级正确。(c) $E(a) =
U/[a\ln(b/a)] = 1000/(5 \times 10^{-4} \times 1.61) = \qty{1.2}{MV/m}$：低于
聚乙烯的 \qty{20}{MV/m}，是安全的；在空气中则将处于临界状态。
\end{solution}

\begin{solution}{exo:b1:potential-capacitors:11}
(a) $W = Q^2/2C = Q^2/8\pi\varepsilon_0R$。(b) $R = e^2/8\pi\varepsilon_0m_ec^2 = 2.31 \times
10^{-28}/(2 \times 8.19 \times 10^{-14}) = \qty{1.4e-15}{m}$（通常所说的“经典
半径”，$e^2/4\pi\varepsilon_0m_ec^2 = \qty{2.8}{fm}$ 则略去了 $\tfrac12$）。(c) $W =
Q^2/2C = Q^2d/2\varepsilon_0S$ 加倍：$\Delta W = Q^2d/2\varepsilon_0S = F\,d$，其中 $F =
Q^2/2\varepsilon_0S = \sigma^2S/2\varepsilon_0$，与直接求得的结果一致。(d) 在 $U$ 固定时，$W =
\varepsilon_0SU^2/2d$ 减半：$\Delta W = -W/2$；外力仍做功 $W/2$
（对 $F = \varepsilon_0SU^2/2x^2$ 从 $d$ 积分至 $2d$）；两部分合计 $W$，
回流入电池，此时电荷 $CU/2$ 以
\omterm{def:b1:dc-circuits:voltage}{电压} $U$ 返回电池。
\end{solution}

\begin{solution}{exo:b1:potential-capacitors:12}
(a) 半径为 $r$ 的球带电 $q = Qr^3/R^3$，表面电势为
$q/4\pi\varepsilon_0r$；添加球壳 $\dd q = 3Qr^2\,\dd r/R^3$ 需做功 $\dd W =
q\,\dd q/4\pi\varepsilon_0r = 3Q^2r^4\,\dd r/4\pi\varepsilon_0R^6$；积分得 $W =
3Q^2/20\pi\varepsilon_0R$。(b) $Q = 92e = \qty{1.47e-17}{C}$：$W = 3 \times 2.17 \times
10^{-34}/(20\pi \times 8.85 \times 10^{-12} \times 7.4 \times 10^{-15}) = \qty{1.6e-10}{J}
= \qty{990}{MeV}$。(c) 每个碎片：电荷 $Q/2$，半径 $R/2^{1/3}$，
能量 $\tfrac14 2^{1/3}W$；两个合计 $2^{1/3}W/2 = 0.63W$：释放
$0.37 \times 990 = \qty{370}{MeV}$。(d) 原子核的“表面张力”：
两个碎片比一个原子核多出表面，这约需
\qty{170}{MeV}；能量平衡接近 \qty{200}{MeV}，主要以
两碎片在相互斥力作用下分离飞开的\omterm{thm:b1:work-and-energy:ke}{动能}
形式出现。
\end{solution}

\begin{solution}{pb:b1:potential-capacitors:1}
\textbf{1.} 重力 $\tfrac43\pi r^3\rho g$ 向下，阻力 $6\pi\eta rv$ 向上，
浮力 $\tfrac43\pi r^3\rho_{\text{空气}}g$ 向上——为 $1.2/900 = 0.13\%$ 的
重力，可忽略不计。\omterm{prop:b1:newton-dynamics:terminal}{终端速度}：$\tfrac43\pi r^3\rho g = 6\pi\eta rv_1$，$v_1 =
2\rho gr^2/9\eta$.

\textbf{2.} $v_1 = \qty{1.0e-4}{m/s}$：$r = \sqrt{9\eta v_1/2\rho g} = \sqrt{9 \times
1.8 \times 10^{-5} \times 10^{-4}/(2 \times 900 \times 9.8)} = \qty{9.6e-7}{m}$，约为
微米量级；$m = \tfrac43\pi r^3\rho = \qty{3.3e-15}{kg}$。

\textbf{3.} $\tau = m/6\pi\eta r = 3.3 \times 10^{-15}/3.3 \times 10^{-10} = \qty{1e-5}{s}$；
距离 $\sim v_1\tau = \qty{1}{nm}$。不可能。

\textbf{4.} $\mathrm{Re} = 1.2 \times 10^{-4} \times 9.6 \times 10^{-7}/1.8 \times 10^{-5} =
6 \times 10^{-6} \ll 1$：斯托克斯定律成立。

\textbf{5.} $eE = 1.6 \times 10^{-19} \times 3 \times 10^5 = \qty{4.8e-14}{N}$，与
$mg = \qty{3.2e-14}{N}$ 相当，仅微米量级的油滴如此。直径为 \qty{10}{\micro m}
的油滴重一千倍；一个\omterm{def:b1:electrostatics-gauss:charge}{元电荷}引起的
速率变化仅为 $0.1\%$，不可察觉。

\textbf{6.} 布朗运动：空气分子的撞击使油滴随机地抖动；
计时结果因此散布，限制了测量精度，
并为有用的油滴尺寸设定了下限。

\textbf{7.} $E = U/d = 5000/0.016 = \qty{3.1e5}{V/m}$；极板远宽于
$d$，因此两平面的叠加给出匀强电场；
等势面为水平面，$V = Uz/d$。

\textbf{8.} $C = \varepsilon_0\pi(0.1)^2/0.016 = \qty{17}{pF}$；$Q = CU = \qty{87}{nC}$；
$W = \tfrac12CU^2 = \qty{0.22}{mJ}$。

\textbf{9.} 上极板为正：$\vect E$ 向下；
受力须向上，故 $q < 0$。$|q| = mgd/U_0 = 3.3 \times 10^{-15} \times 9.8 \times 0.016/
1080 = \qty{4.8e-19}{C} = 3e$：$q = -3e$。

\textbf{10.} 匀速上升：$|q|E = mg + 6\pi\eta rv_2$，且 $mg =
6\pi\eta rv_1$，故 $|q|E = 6\pi\eta r(v_1 + v_2)$，$|q| = 6\pi\eta r(v_1 +
v_2)d/U$。$v_2 = (|q|E - mg)/6\pi\eta r = (1.5 \times 10^{-13} - 3.2 \times 10^{-14})/
3.3 \times 10^{-10} = \qty{3.6e-4}{m/s}$：\qty{1}{mm} 历时 \qty{2.8}{s}。

\textbf{11.} 平衡判断是对一个漂移、抖动的微粒作零值判断，
其电荷在此期间可能改变；沿固定距离的计时可随意重复，
$r$ 一次性由 $v_1$ 确定，同一油滴便
可得到一系列电荷值。

\textbf{12.} $|q|U = 3 \times 1.6 \times 10^{-19} \times 5000 = \qty{2.4e-15}{J}$
（\qty{15}{keV}）为跨越间隙的功；在 \qty{1}{cm} \omterm{def:b1:dc-circuits:sources}{内阻}力做功：$(|q|E - mg) \times 0.01 =
\qty{1.2e-15}{J}$，转化为空气中的热量。

\textbf{13.} 除以 $1.60$：$3.00$、$5.01$、$4.01$、$1.99$、$6.02$——
整数 $3, 5, 4, 2, 6$。

\textbf{14.} $q/n$：$1.600, 1.602, 1.605, 1.595, 1.605$；均值 $1.601$，
标准差 $0.004$，均值的不确定度 $0.004/\sqrt5 =
0.002$：$e = (1.601 \pm 0.002) \times 10^{-19}$~C。

\textbf{15.} X 射线使空气电离；油滴每次俘获一个离子，
每个离子所带电荷为 $e$ 的整数倍。

\textbf{16.} $r \propto \eta^{1/2}$（由 $v_1$ 可得），$m \propto r^3 \propto \eta^{3/2}$，且
$|q| = mgd/U_0 \propto \eta^{3/2}$：$1\%$ 的 $\eta$ 误差导致 $1.5\%$ 的 $e$ 误差。密立根所用的
$\eta$ 偏低 $0.4\%$，因此他的 $e$ 偏低 $0.6\%$。

\textbf{17.} 夸克从不单独存在：它们被束缚在强子之内，
而每个自由物体的电荷都是\omterm{def:b1:electrostatics-gauss:charge}{元电荷}的整数倍。

\textbf{18.} $N_A = F/e = 96485/1.602 \times 10^{-19} = \qty{6.02e23}{mol^{-1}}$。

\textbf{19.} $m_e = 1.602 \times 10^{-19}/1.76 \times 10^{11} = \qty{9.1e-31}{kg}$：比氢原子轻
$1840$ 倍。

\textbf{20.} 极板间为水平面，在边缘附近向外弯曲；
竖直电\omterm{def:b1:electrostatics-gauss:field}{场线}在边缘向外展开，展开的宽度约为 $d = \qty{1.6}{cm}$ 量级，
相对于 \qty{20}{cm} 的直径很小，
而油滴是在中心处观察的。

\textbf{21.} $E_p = qUz/d + mgz$，关于 $z$ 呈线性：没有极小值；
静止油滴要求斜率为零，即 $q = -mgd/U$，此时平衡为随遇平衡，
任何残余漂移都会持续，
这正是判断平衡困难的原因。

\textbf{22.} $2.4 \times 10^{-15}$ 与 $2.2 \times 10^{-4}$~J 相比：$10^{-11}$。油滴
不扰动电场。

\textbf{23.} 电场向上，负油滴所受电场力向下，油滴以 $(mg
+ |q|E)/6\pi\eta r = (3.2 \times 10^{-14} + 1.5 \times 10^{-13})/3.3 \times 10^{-10} =
\qty{5.6e-4}{m/s}$ 的速度下落。

\textbf{24.} $F = 10^4 \times 1.6 \times 10^{-19} \times 3.1 \times 10^5 = \qty{5e-10}{N}$
与 $\qty{9.8e-9}{N}$ 相比：为重力的 $5\%$，无法托起；
且若有 $10^4$ 个电荷，改变一个 $e$ 是不可察觉的。
油雾所产生的油滴很小、不易蒸发，且为球形。

\textbf{25.} 测得：$v_1$、$v_2$ 与 $U$（$d$、$\eta$、$\rho$ 已知）；
由此算出 $r$、$m$，进而得到 $q$，发现其总是
$e = (1.601 \pm 0.002) \times 10^{-19}$~C（$0.1\%$）的整数倍；由此求得 $N_A = F/e$ 与
$m_e = e/(e/m_e)$。
\end{solution}
